\documentclass[14pt,a4paper]{extarticle}
\usepackage[left=3cm,right=2cm,top=2cm,bottom=2cm,headheight=18pt]{geometry}
\usepackage{fontspec}
\setmainfont{Liberation Serif}
\setsansfont{Liberation Serif}
\setmonofont{Liberation Mono}
\usepackage{unicode-math}
\setmathfont{texgyretermes-math.otf}
\usepackage{amsmath,graphicx,booktabs,longtable,array,enumitem}
\usepackage{titlesec,fancyhdr}
\usepackage[hidelinks,unicode]{hyperref}
\usepackage{url}
\urlstyle{same}
\setlength{\parindent}{1.27cm}
\setlength{\parskip}{0pt}
\linespread{1}
\setlength{\emergencystretch}{2em}
\clubpenalty=10000
\widowpenalty=10000
\titleformat{\section}{\normalfont\bfseries}{\thesection.}{0.5em}{}
\titleformat{\subsection}{\normalfont\bfseries}{\thesubsection.}{0.5em}{}
\titlespacing*{\section}{0pt}{1.2em}{0.5em}
\titlespacing*{\subsection}{0pt}{0.9em}{0.3em}
\setlist{nosep,leftmargin=1.27cm}
\pagestyle{fancy}
\fancyhf{}
\renewcommand{\headrulewidth}{0pt}
\renewcommand{\footrulewidth}{0pt}
\fancyhead[C]{\ifnum\value{page}>1\thepage\fi}
% Liberation Serif is a disclosed Times-compatible substitute, not Times New Roman.
% Current statutory geometry/size baseline: O‘RQ-682 appendix §§43–51.
\hyphenpenalty=10000
\exhyphenpenalty=10000
\pretolerance=10000
\tolerance=10000
\setlength{\emergencystretch}{4em}
\renewcommand{\normalsize}{\fontsize{14bp}{16.8bp}\selectfont}
\normalsize
\begin{document}
\noindent Muallif: Abduxoliq Akmal ugli Ashuraliyev. Tahrir sanasi: 2026-yil 1-oktabr.

\begin{flushright}Loyiha\end{flushright}
\begin{center}\textbf{NORMATIV-HUQUQIY HUJJAT LOYIHASIGA\\ KIRITISH UCHUN TAKLIF}\par
\medskip\textbf{Oilalarni ko‘p o‘lchovli baholash hisob-kitoblarini tekshirish tartibi to‘g‘risida}\end{center}
\noindent PF-193-son Farmonning 14-bandi bo‘yicha tayyorlanayotgan hujjatga kiritish uchun mualliflik taklifi. Asosiy ruscha murojaat matniga mos o‘zbekcha tarjima. Ushbu matn rasmiy loyihani davlat tilida tayyorlash uchun ham taqdim etiladi; inglizcha tarjima ilova qilinadi. Yakuniy hujjat turi va tasdiqlash rekvizitlari ishlab chiquvchi tomonidan belgilanadi.\par
\noindent\hspace*{1.27cm}1. Oilalarni Ijtimoiy reyestrga kiritish, toifalarga ajratish va tegishli xizmatlarga yo‘naltirishda qo‘llaniladigan ko‘p o‘lchovli baholash uslubiyotining hisob-kitoblarini tekshirish tartibi ilovaga muvofiq tasdiqlansin.\par
\noindent\hspace*{1.27cm}2. O‘zbekiston Respublikasi Prezidenti huzuridagi Ijtimoiy himoya milliy agentligi (keyingi o‘rinlarda — Agentlik) O‘zbekiston Respublikasi Iqtisodiyot va moliya vazirligi (keyingi o‘rinlarda — Vazirlik) bilan birgalikda ko‘p o‘lchovli baholash uslubiyotini belgilangan tartibda tasdiqlash uchun tayyorlashda mezonlar, hisoblash qoidalari, toifalarga ajratish chegaralari hamda ularning huquqiy va ilmiy asoslarini yagona tavsifda aks ettirsin. Dasturiy vositaning texnik sozlamalari orqali mustaqil ravishda yangi mezon yoki rad etish asosi joriy etilishiga yo‘l qo‘yilmasin.\par
\noindent\hspace*{1.27cm}3. Ijtimoiy himoya milliy agentligi hisob-kitoblarning tasdiqlangan uslubiyotga muvofiqligi tekshirilishini, natijalarni qayta hisoblash imkoniyatini va qaror asoslarini saqlashni ta’minlasin. Tekshiruv oilalarning murojaatlarini ko‘rib chiqishning qonunchilikda belgilangan muddatlarini uzaytirish uchun asos bo‘lmaydi.\par
\noindent\hspace*{1.27cm}4. Agentlik Vazirlik bilan birgalikda “Normativ-huquqiy hujjatlar to‘g‘risida”gi Qonunning 22-moddasi asosida uslubiyot va tekshiruv materiallarini tayyorlashni tashkil etsin. Zarur bilim va tajribaga ega fuqarolar roziligi bilan ishchi guruhga kiritilishi, ilmiy tavsiyalar va mutaxassislar maslahati olinishi mumkin. Agentlik uslubiy, tizim va yuridik vazifalarni taqsimlasin hamda mezon pasportlari, manbalarni moslashtirish, sinovlar ro‘yxati, bayonnomalar va nuqson qaydlari uchun ichki texnik reglamentni tasdiqlasin. Reglament ushbu Tartibni bajaradi, muhtojlik mezonlari, rad etish asoslari, arizachi majburiyatlari yoki qonuniy muddatlar o‘zgarishini joriy etmaydi. Loyihada qatnashish mezonlarni tasdiqlash yoki himoyalangan yozuvlarni olish vakolatini bermaydi.\par
\noindent\hspace*{1.27cm}5. Agentlik va Vazirlik hisoblash tavsifini reyestrdagi baholash, toifa, yordam, shikoyat va axborot almashish qoidalari bilan asosiy loyihada tegishli qoida birinchi qo‘llanishidan oldin muvofiqlashtirsin. Yuqori hujjatdagi ziddiyat vakolatli qabul qiluvchi organga kiritiladi. Texnik tuzatish qaror sanasida qonuniy qo‘llanadigan qoidani bajaradi; almashtirilgan mezonni uzaytirmaydi, joriy etish sanasini o‘zgartirmaydi va xodimga qonuniy tartibdan tashqarida toifa qarorini qo‘lda kiritish vakolatini bermaydi.\par
\noindent\hspace*{1.27cm}6. Tekshiruv mavjud tizimni ishlab chiqish, texnik kelishish va o‘zgarishlarni boshqarish jarayoniga kiritilsin. Teng dalilni ta’minlaydigan mavjud spetsifikatsiya, sinov, qayd va qonuniy nazorat ishlatilsin; takroriy nazorat va xarid talab qilinmasin. Qo‘shimcha xarajat majburiyatidan oldin Agentlik va Vazirlik qo‘shimcha ish, ajratilgan xodim vaqti, shartnoma qamrovi va qonuniy mablag‘ni belgilasin. Ishlov o‘zgarishi unga tegishli yuklama va vaqt qaydlari bilan, zarur bo‘lsa ruxsat etilgan muhitdagi quvvat sinovi bilan asoslanadi. Qiymat formulasi, shtat jadvali va umumiy budjet dalilni almashtirmaydi. Tartib alohida respublika miqyosidagi qo‘lda ko‘rib chiqish xizmati, parallel tizim, arizachi to‘lovi yoki qo‘shimcha ariza bosqichini yaratmaydi.\par
\noindent\hspace*{1.27cm}7. Huquqiy oqibatli qo‘llashdan oldin Agentlik qonuniy qoida, zarur texnik kelishuvlar, nazorat natijalari va mavjud izoh hamda e’tiroz yo‘llari tayyorligini 25-banddagi chiqarilish qaydida tasdiqlasin. Muvaffaqiyatsiz tekshiruv 26-band bo‘yicha hal qilinadi. Farovonlik empirik tadqiqoti va alohida ochiq tekshiruv hisoboti joriy etish sharti emas. Texnik qabul qonuniy kelishuvlarni almashtirmaydi, 2027-yil 1-yanvar boshlanishini o‘zgartirmaydi va arizachi muddatini uzaytirmaydi.\par
\clearpage\setcounter{page}{1}
\begin{flushright}\fontsize{12bp}{14.4bp}\selectfont 1-materialga ilova\\Mualliflik taklifi\\2026-yil 1-oktabrdagi tahrir\end{flushright}
\begin{center}\textbf{Ko‘p o‘lchovli baholash hisob-kitoblarini tekshirish tartibi}\end{center}
\subsection*{1-bob. Umumiy qoidalar}
\noindent\hspace*{1.27cm}1. Tartib hisoblashning tasdiqlangan uslubiyotga muvofiqligi, ishlatilgan ma’lumotlarning kuzatilishi va qaror asoslarini saqlashga taalluqli. Nazorat misollari joriy etish va tegishli o‘zgarishdan oldin mavjud chiqarilish jarayonida tekshiriladi; har arizaga ikkinchi qo‘lda hisoblash yuklanmaydi. Sezgirlik yoki farovonlik tadqiqoti alohida kelishilgan savol, qonuniy kirish va resursni talab qiladi; yordam olish yoki tasdiqlangan qoidani ishga tushirishning qo‘shimcha sharti emas. Xizmatga yo‘naltirish shu qoidadan kelib chiqadigan qismda tekshiriladi; tarmoq xizmat shartlari saqlanadi.\par
\noindent\hspace*{1.27cm}2. Ushbu Tartibda “Yagona milliy ijtimoiy himoya” axborot tizimi (keyingi o‘rinlarda — YAMIH) va Agentlikning “Inson” ijtimoiy xizmatlar markazlari (keyingi o‘rinlarda — “Inson” markazlari) qisqartirilgan nomlari qo‘llaniladi. Quyidagi tushunchalardan foydalaniladi:\par
\textbf{uslubiyot tavsifi} — baholashning mezonlari, ma’lumot manbalari, hisoblash qoidalari, toifalarga ajratish shartlari va qo‘llanish muddatlarini qayd etuvchi hujjat;

\textbf{hisoblash talqini} — hisoblashda qo‘llangan tasdiqlangan uslubiyot identifikatori, dasturiy chiqarilish identifikatori va amal qilish sanalari belgilangan parametrlar sozlamasini bog‘lovchi qayd;

\textbf{tekshiruv bayonnomasi} — tekshiruv shartlari, natijalari, aniqlangan nomuvofiqliklar va ular bo‘yicha ko‘rilgan choralar qaydi;

\textbf{hal etilmagan ma’lumot} — kelib tushmagan, tasdiqlanmagan, o‘zaro zid yoki uslubiyotda belgilangan yangilanish talabiga javob bermaydigan ma’lumot;

\textbf{nazorat hisob-kitobi} — tasdiqlangan qoidalar asosida mustaqil qayta hisoblash natijasi.

\noindent\hspace*{1.27cm}3. Hisoblash uslubiyoti oila uchun yordam yoki xizmat tayinlash vakolatini almashtirmaydi. Toifa, muayyan yordamga huquq va yordam miqdori qonunchilikda belgilangan tartibda alohida aniqlanadi.\par
\subsection*{2-bob. Uslubiyot tavsifi va hisoblash talqinlari}
\noindent\hspace*{1.27cm}4. Har bir mezon tavsifida huquqiy asos, maqsad, vakolatli manba, birlik, hisobga olinadigan davr, yangilash talabi va hisoblash qoidasi belgilanadi. Amal sanalari, o‘zgaruvchan parametrlar, takroriy yoki ustma-ust axborot va istisnolar tasdiqlangan uslubiyotga muvofiq aniqlanadi. Agentlik ichki texnik reglamenti pasport maydonlari, manbalarni moslashtirish va aniqlik tekshiruvlarini belgilaydi; shakllar o‘zicha yangi qoida yoki rad etish asosini yaratmaydi.\par
\noindent\hspace*{1.27cm}5. Tavsif hisoblash formulalari, shartli qoidalar, koeffitsiyentlar, mezonlar o‘rtasidagi bog‘liqliklar, istisnolar, yaxlitlash tartibi va toifalarga ajratish chegaralarini qamrab oladi. Chegaraga aynan teng qiymat qaysi toifaga tegishli ekanligi aniq belgilanadi. Qonunchilikdagi istisnolar va ustuvorlik qoidalari qo‘llangandan keyin bir holat uchun bir-biriga zid toifalar belgilanmasligi, toifaga kiritish va rad etish shartlari o‘rtasida yechimsiz holat qolmasligi tekshiriladi. Aniqlashtirish talab qilinadigan holat toifaga kiritish yoki rad etish qarori bilan aralashtirilmaydi. Dasturiy vosita tavsifda nazarda tutilmagan baholash qoidasini qo‘llamaydi.\par
\noindent\hspace*{1.27cm}6. Har bir hisoblash talqini uchun tasdiqlash asosi, qo‘llanish boshlanishi va tugashi, oldingi talqindan farqlari hamda tekshiruv bayonnomasi saqlanadi. Huquqiy oqibatga ta’sir qiladigan o‘zgarish vakolatli organ tomonidan belgilangan tartibda tasdiqlanadi. Huquqiy uslubiyot, dasturiy chiqarilish va parametrlar sozlamasi alohida aniqlanadigan qaydlar bilan hisoblash talqiniga bog‘lanadi; huquqiy amal sanalari texnik joriy etish sanalaridan ajratiladi. O‘zgarmagan qonuniy qoidani bajarishga qaratilgan dasturiy tuzatish yangi siyosat mezoni deb ko‘rsatilmaydi. Sozlamada manbalarni maydonlarga moslashtirish, qiymatlarni aylantirish va yaxlitlash tartibi qayd etiladi; natijaga ta’sir qilishi mumkin bo‘lgan o‘zgarishlar tekshiriladi.\par
\noindent\hspace*{1.27cm}7. Elektron tavsif uslubiyotni tasdiqlagan e’lon qilingan hujjat, amal sanalari va hisoblash talqinini ko‘rsatadi. Qoidaning ochiq qismi Agentlikning mavjud rasmiy resursida ochiladi; vakolatli e’lon qilingan tavsifni ikkinchi hisobotga ko‘chirish o‘rniga havola berish mumkin. Cheklangan qism faqat vakolatli tekshiruvchi va shikoyat ko‘ruvchiga taqdim etiladi. Dastur kodi, himoyalangan yozuv va infratuzilma yoki xavfsizlik tafsilotini nashr qilish talab etilmaydi; mustaqil mezon yoki istisno yaratilmaydi. Qonuniy vaqtinchalik uzilish qoidasi 11-banddagi talqin qaydida belgilanadi.\par
\subsection*{3-bob. Ma’lumotlar sifati va uzilishlar}
\noindent\hspace*{1.27cm}8. Har bir hisoblashda foydalanilgan ma’lumotning manbasi, tegishli davri, olingan vaqti va tekshirilganlik holati qayd etiladi. Haqiqiy nol qiymat hal etilmagan ma’lumotdan ajratiladi. Manba yozuvi ifodalaydigan sana yoki davr uning olingan vaqtidan ajratiladi. Zarur kirish qiymatlari yoki qonuniy saqlangan manba nusxasi 28-bandga rioya qilgan holda hisoblashni qayta tiklash imkonini berishi kerak. Tasdiqlangan qiymat, olinmagan ma’lumot, yangilash talab qilinishi, ziddiyatli ma’lumot va noto‘g‘ri format alohida saqlanadi. Ishlatilgan yozuvning manba hamda yozuv identifikatori, tegishli davri, olish sanasi va qo‘llangan o‘zgartirish qayd etiladi. Ziddiyatda ikkala qiymat va tanlov asosi saqlanadi; keyinroq olinganlik o‘zicha ko‘proq ishonchlilikni bildirmaydi. Xodim tuzatishi muallif, vaqt va qonuniy asosni qamraydi.\par
\noindent\hspace*{1.27cm}9. Hal etilmagan ma’lumotni oilaning holatini yomonlashtiruvchi qiymatga avtomatik tenglashtirishga yo‘l qo‘yilmaydi. Qonunchilikda bevosita nazarda tutilgan normativ qiymat qo‘llansa, uning huquqiy asosi va qo‘llanish sababi hisob-kitobda alohida ko‘rsatiladi.\par
\noindent\hspace*{1.27cm}10. Hal etilmagan ma’lumot sababli hisoblash imkonsiz bo‘lsa, “Inson”ning vakolatli xodimi masalani YAMIHda qayd etib, mas’ul manba egasi yoki operatoriga yuboradi. Dastur nuqsoni tizim bo‘linmasiga yuboriladi; huquqiy qoida noaniqligi uslubiy bo‘linma va yuridik xizmatda ko‘rilib, rasmiy sharh yoki o‘zgartirish zarur bo‘lsa vakolatli organga kiritiladi. Agentlik xodimlari uning vakolatini almashtirmaydi. Mas’ul xodim qonuniy muddatni kuzatadi, oilaga natija va shikoyat yo‘lini bildiradi, hal etilmagan masalani muddat tugashidan oldin rahbarga chiqaradi. Holat yangi rad etish yoki tugatish asosi emas. Davlat tizimi ma’lumoti o‘rniga oila hujjati faqat qonunchilik ruxsat etgan holda talab qilinadi.\par
\noindent\hspace*{1.27cm}11. Texnik nosozlik, axborot almashishdagi uzilish yoki elektron ma’lumot to‘liq olinmagani munosabati bilan qonunchilikda nazarda tutilgan vaqtinchalik algoritm o‘zgarishi joriy etilganda, vakolatli mansabdor shaxs uning asosini, qo‘llanish doirasini, boshlanish vaqtini, vaqtinchalik hisoblash qoidalarini va tiklash shartlarini qayd etadi. Vaqtinchalik talqin qo‘llangan hisob-kitoblar alohida belgilab boriladi. Tizim tiklangach, ularning tekshiruvi va zarur qayta hisoblash qonunchilikda belgilangan tartibda amalga oshiriladi. Vaqtinchalik qayd ta’sirlangan manba, hisoblash va talqinlarga bog‘lanadi; qaror qilgan xodim, uning huquqiy vakolati va rejim tugash sanasi ko‘rsatiladi. Tiklangach vaqtinchalik va tiklangan hisoblar solishtiriladi, farq va vakolatli organning individual qarorlari saqlanadi. Uzilishning yopilishi manba mavjudligi va zarur sinovlar tugashi bilan tasdiqlanadi, vaqtinchalik qaydni o‘chirish bilan emas.\par
\subsection*{4-bob. Hisob-kitoblarni tekshirish}
\noindent\hspace*{1.27cm}12. Joriy etish yoki ball, toifa yoxud qaror asosiga ta’sirli o‘zgarishdan oldin dastur tasdiqlangan huquqiy qoidadan mustaqil asoslangan nazorat javoblariga qiyoslanadi. Mavjud qabul va texnik muvofiqlik sinovlari tegishli chegara, istisno, ularning o‘zaro ta’siri, ma’lumot sifati holati va amal sanalarini qamraydi. Hujjatlangan teng sinov faqat ushbu Tartib uchun takrorlanmay, qayta ishlatiladi. Kutilgan javob faqat sinalayotgan dasturdan olinmaydi. Ichki reglament misol rekvizitlari va rejalangan hamda bajarilgan tarkibni solishtirishni belgilaydi. Qo‘llanadigan yo‘l tushirib qoldirilishi, bo‘sh taqqoslash tarkibi yoki hal etilmagan qoida ziddiyati muvaffaqiyatli tekshiruv deb ko‘rsatilmaydi.\par
\noindent\hspace*{1.27cm}13. Tekshiruvning maqsadi, qamrab olinadigan davri va holatlari, tanlash qoidalari, solishtirish mezonlari hamda natijani baholash shartlari natijalar olinmasidan oldin belgilanadi. Keyingi o‘zgartirishlar sababi bilan qayd etiladi. Sintetik sinov natijalari haqiqiy oilalar ehtiyojini to‘g‘ri aniqlash isboti sifatida qo‘llanmaydi.\par
\noindent\hspace*{1.27cm}14. Ball va toifa bo‘yicha nomuvofiqliklar alohida hisoblanadi. Ma’lumotlar yoki uslubiyotdagi tekshirilayotgan o‘zgarishlar natijasida toifasi almashadigan holatlar soni, tegishli jami holatlar soni va o‘zgarishning mazmuni ko‘rsatiladi. Hal etilmagan holatlar hisobotdan yashirilmaydi va aniq toifa belgilangan holatlardan ajratiladi. Toifa nomuvofiqligi ko‘rsatkichi maxraji bir xil tasdiqlangan qoida va aynan bir xil kirishlar bo‘yicha ikkala hisoblashda ham toifasi aniq belgilangan qiyoslanadigan holatlarnigina qamrab oladi. Sezgirlik ko‘rsatkichi maxraji belgilangan o‘zgarishdan oldin va keyin ham toifasi aniq bo‘lgan juft holatlarnigina qamrab oladi. Har bir ko‘rsatkich uchun surat, maxraj, chiqarib tashlangan va hal etilmagan holatlar soni ochiqlanadi; maxraj nol bo‘lsa, ko‘rsatkich nol emas, aniqlanmagan deb qayd etiladi.\par
\noindent\hspace*{1.27cm}15. Ehtiyojni noto‘g‘ri aniqlash ko‘rsatkichlari mustaqil asoslangan taqqoslash ma’lumotlari mavjud bo‘lgandagina hisoblanadi. Oldingi avtomatik qaror, eski toifa yoki shikoyat natijasi o‘z-o‘zidan to‘g‘ri ehtiyoj bahosi hisoblanmaydi. Shikoyat natijalari shikoyat qilmagan oilalarni ham qamrab oluvchi xatolik ko‘rsatkichi sifatida talqin etilmaydi.\par
\noindent\hspace*{1.27cm}16. Bayonnomada tasdiqlangan qoida va tekshirilgan amalga oshirish, tekshiruv qamrovi va asosi, ijrochilar, natijalar, tekshirilmagan yoki hal etilmagan masalalar, aniqlangan oqibatlar va tuzatish qarorlari ko‘rsatiladi. Qoidaga moslik, ma’lumot sifati va qaror asoslari to‘liqligi haqiqiy muhtojlikni empirik aniqlashdan ajratiladi. Yakunlar teng bo‘lsa ham oraliq nomuvofiqliklar saqlanadi. Dasturlar mosligi o‘zicha qonuniy kutilgan javobni belgilamaydi. Ichki texnik reglament dalil va qayd shakllarini belgilaydi; himoyalangan kirishlar qonuniy foydalaniladigan qismda qoladi.\par
\noindent\hspace*{1.27cm}17. Tasdiqlangan uslubiyotga zid hisoblash aniqlanganda, tegishli dasturiy talqin tuzatiladi va takroran tekshiriladi. Ta’sirlangan holatlar aniqlanib, vakolatli organ tomonidan qonunchilikka muvofiq qayta ko‘rib chiqiladi. Tekshirilmagan talqin muvofiqligi tasdiqlangan talqin sifatida qo‘llanmaydi; bu qoida murojaatlarni ko‘rib chiqish yoki qonuniy yordamni ko‘rsatishni to‘xtatishga asos bo‘lmaydi. Nuqson qaydi qoida va talqin, takrorlanadigan misol, ta’sirlangan qarorlar, mas’ul, tuzatish usuli, tatbiq etiladigan muddat va takroriy tekshiruv natijasini qamraydi. Qarorlar faqat nuqsonni ko‘rsatgan shikoyatga emas, qo‘llangan talqin, sana va shartlarga ko‘ra tanlanadi. O‘zgargan komponent va bog‘liq hisoblash yo‘llari qayta sinaladi; ta’sir sohasini aniqlab bo‘lmasa, barcha tegishli hisoblash zanjiri tekshiriladi. Yopilish sinov dalili va individual qarorlar zarur qayta ko‘rilishi qaydi bilan tasdiqlanadi.\par
\subsection*{5-bob. Qaror asoslari va ularni qayta ko‘rib chiqish}
\noindent\hspace*{1.27cm}18. Oila bo‘yicha hisoblash natijasi bilan birga qo‘llangan talqin, tegishli mezonlar, ma’lumot manbalari va davrlari, ball, toifani belgilash qoidasi, normativ qiymatlar hamda vakolatli xodim kiritgan tuzatishning asosi saqlanadi. Qayta hisoblashda oldingi qayd o‘chirib yuborilmaydi. Qayd zarur oraliq hissalar, kiritish yoki chiqarish sabablari, aylantirish va yaxlitlash bosqichlari hamda 6-banddagi identifikatorlarni saqlaydi. Bir-birini qoplaydigan tarkibiy tafovutlar jami natija bir xil bo‘lgani sababli yashirilmaydi.\par
\noindent\hspace*{1.27cm}19. Ariza beruvchiga qaror natijasi bilan bir vaqtda uning huquqiy va faktik asoslari, hisoblash talqini, asosiy qiymatlar, davrlar, normativ qiymat yoki istisno qo‘llanish sababi hamda e’tiroz bildirish tartibi elektron shaxsiy kabinet orqali taqdim etiladi. Elektron kabinetdan foydalanmaydigan shaxsga “Inson” markazi orqali mazkur axborotni olish imkoniyati ta’minlanadi. Uchinchi shaxslarning himoyalangan ma’lumotlari oshkor etilmaydi. Natija haqida SMS-xabar yuborish ushbu tushuntirish va kirish imkoniyatini almashtirmaydi. Izoh qaror raqami va sanasi, tatbiq etilgan norma, arizachining tegishli qiymatlari va davrlari, asosiy o‘zgartirishlar, toifa va muayyan yordam yoki yo‘naltirish shartlari, vakolatli e’tiroz oluvchi va qonuniy muddatni beradi. Normativ daromad tasdiqlangan haqiqiy daromaddan ajratiladi. Faqat avtomatik ishlovga asoslangan qaror uchun oldindan qonun ruxsat bergan asos, ishlov va mumkin bo‘lgan huquqiy oqibatni tushuntirish tartibi qayd etiladi. Ayrim ma’lumotning qonuniy cheklovi qolgan ochiq asoslarni taqdim etishga to‘sqinlik qilmaydi.\par
\noindent\hspace*{1.27cm}20. Faqat avtomatlashtirilgan ishlov berish asosidagi qarorga e’tiroz ma’lumotlar mulkdori yoki operatori tomonidan ko‘rib chiqiladi hamda natijasi haqida subyektga o‘n kun ichida yozma shaklda xabar beriladi. Bunday e’tiroz YAMIH axborot tizimi yoki “Inson” markazi orqali mulkdor yoki operatorga yuborilishi mumkin; axborot tizimida qayd etilishi muddat hisoblanishini kechiktirmaydi. Ma’lumotni tuzatish talabi uning vakolatli manbasiga yuboriladi. Subyekt murojaati bo‘yicha shaxsga doir ma’lumotlar mulkdor yoki operator tomonidan murojaat berilgan paytdan e’tiboran uch kundan kechiktirmay o‘zgartiriladi yoki to‘ldiriladi; haqiqatga mos emaslik aniqlansa, tuzatish darhol amalga oshiriladi. Apellyatsiya kengashiga shikoyat kengashning qonunchilikda belgilangan predmeti doirasida ko‘riladi; kengash vakolatiga kirmaydigan hisoblash e’tirozi shu sababli ko‘rib chiqilmasdan qoldirilmaydi, balki tegishli vakolatli organga yuboriladi. Sudga yoki boshqa vakolatli organga murojaat qilishning qonuniy tartibi va muddatlari saqlanadi. Ushbu band yangi majburiy sudgacha bosqichni joriy etmaydi.\par
\subsection*{6-bob. Tadqiqot, hisobot va nazorat}
\noindent\hspace*{1.27cm}21. Tadqiqot va ilmiy tekshiruv uchun ma’lumotlardan foydalanishdan oldin ularga ishlov berishning qonuniy asosi, maqsadi, zarur hajmi va himoya choralari belgilanadi. Tadqiqot uchun shaxsga doir ma’lumotlarni egasizlantirish qonunchilik talablariga muvofiq amalga oshiriladi. Identifikatorni almashtirishning o‘zi egasizlantirish deb hisoblanmaydi. Sog‘liq haqidagi ma’lumotlarga doir maxsus talablar bajariladi. Tadqiqotga ruxsat qaydida qonuniy savol, zarur maydon va davr, vakolatli operator, ruxsat etilgan amallar, saqlash muddati va oshkor qilish nazorati ko‘rsatiladi. Fuqaroning tayyorlashdagi ishtiroki oila ma’lumotiga ishlov uchun oilaning roziligi emas. Sog‘liq haqidagi ma’lumot faqat maxsus qonuniy shartlar bilan olinadi. Tadqiqot hisoblash alohida qonuniy individual qarorsiz toifa yoki to‘lovni o‘zgartirmaydi.\par
\noindent\hspace*{1.27cm}22. Tadqiqotning hisob-kitoblari Agentlik nazoratidagi muhitda bajariladi. Tashqi ilmiy ishtirokchiga shaxsga doir ma’lumotlarni olish vakolati ushbu Tartib bilan berilmaydi. Himoyalangan ma’lumotlar, ularni qayta aniqlash imkonini beradigan kichik guruhlar va bog‘langan jadvallar e’lon qilinmaydi. Taqdim etilgan kod bajarilishidan oldin vakolatli operator uning tarkibi, talqini, zarur kirish huquqlari va tashqi ulanishlarini tekshiradi. Ishlov berish eng kam zarur huquqlar bilan bajariladi; himoyalangan kirish qiymatlarini tashqi xizmatlarga uzatish yoki ochiq jurnallarga kiritish taqiqlanadi. Bajarishga ruxsat va oshkor etish nazorati natijalari bayonnomada qayd etiladi. Kodni taqdim etish uni amaldagi tizimda avtomatik bajarishga ruxsat bermaydi.\par
\noindent\hspace*{1.27cm}23. Tekshiruv bayonnomasi mavjud chiqarilish va nazorat qaydlari bilan ichki muhitda saqlanadi. Unda qoida va dastur talqini, sana, soha, bajarilgan va chiqarilgan misol, taqqoslanadigan va hal etilmagan son, isbotlangan tafovut, tuzatish va cheklov beriladi. Teng mavjud qaydga takroriy jurnal yaratmay havola berish mumkin. Axborot mavjud qonuniy hisobot, arizachi kirishi va nazorat tartibida ochiladi. Har chiqarilishdan oldin alohida hisobot, yangi davriy statistik nashr yoki oiladan yangi ma’lumot yig‘ish belgilanmaydi. Qonuniy talab etilgan yoki ruxsat berilgan ochiq ko‘chirma oshkor qilish tekshiruvidan o‘tib, soha va cheklovni beradi; himoyalangan va bog‘lab aniqlanadigan oila yozuvlari chiqariladi.\par
\noindent\hspace*{1.27cm}24. Talqinlarni saqlash, tekshiruv natijalari va tuzatishlar ustidan nazorat amaldagi qonunchilikda belgilangan vakolatlar doirasida tashkil etiladi. Tekshiruv natijalari uslubiyotni takomillashtirish bo‘yicha takliflar tayyorlashda qo‘llaniladi; yordamni avtomatik qisqartirish uchun mustaqil huquqiy asos bo‘lmaydi.\par
\subsection*{7-bob. Talqinni qabul qilish, tuzatish va saqlash}
\noindent\hspace*{1.27cm}25. Agentlik rahbari yoki qonuniy vakolatli o‘rinbosar hisoblash muvofiqligini mavjud chiqarilish qaydida uslubiy, tizim va yuridik xulosalar bilan tasdiqlaydi. Qayd qonun talab qilgan kelishuvlar, amal sanalari, mustaqil nazorat natijalari, qayta ishlatilgan teng sinov, qo‘shimcha ish va resursni bog‘laydi. Ochiq masalalar va 26-banddagi harakat qayd etiladi. Ichki tasdiq vakolatli texnik yoki kiberxavfsizlik kelishuvini almashtirmaydi, muhtojlik siyosati yoki rasmiy sharhlash vakolatini bermaydi. Qarorga ta’sir qiladigan isbotlangan tafovut ahamiyatsiz qolgan vazifa deb hisoblanmaydi.\par
\noindent\hspace*{1.27cm}26. Bir tasdiqlangan qoida, bir xil kirish va tasdiqlangan yaxlitlashda ball, toifa yoki qaror asosi bo‘yicha isbotlangan tafovutli chiqarilish tuzatiladi va ta’sirlangan yo‘llar joriy etishdan oldin qayta sinaladi. Ta’sirlanmagan yo‘l faqat ajratilishi va muvofiqligi isbotlansa davom etadi. Oldingi dasturiy chiqarilish faqat qaror sanasida qonuniy qo‘llanadigan qoidani bajarib, tegishli sinovlari va ishlashga yaroqliligi tekshirilgan bo‘lsa saqlanadi yoki tiklanadi; chiqarilishning eski ekanligi bekor qilingan uslubiyotni uzaytirmaydi. Ishlashda nuqson topilsa, tizim bo‘linmasi soha va ta’sirlangan arizalarni qayd etib, ustuvor tuzatadi va individual qarorlarni qonuniy qayta ko‘rishga yuboradi. Kelib tushish sanasi, qonuniy muddat va vakolatli qaror tartibi saqlanadi; texnik kutish yangi rad asosi yoki qonuniy yordamni to‘xtatish vakolati emas. Axborot almashish uzilishi vakolati faqat o‘z qonuniy predmeti doirasida qo‘llanadi. Toifani qo‘lda kiritish, almashtirilgan baholash qoidasini qo‘llash, qarzdan ozod qilish yoki yanvar boshlanishini to‘xtatish vakolati berilmaydi. Yo‘q yoki zid huquqiy qoida Agentlik va Vazirlik tomonidan vakolatli qabul qiluvchi yoki sharhlovchi organga birinchi qo‘llanishdan oldin hal qilish uchun kiritiladi; ichki sozlama yoki sharh qonunni o‘zgartirmaydi. Texnik zaxira qonuniy qoida yo‘qligini hal qilmaydi. Cheklangan nazorat belgilangan sohani tekshiradi, barcha kirish yoki haqiqiy muhtojlik aniqligini emas.\par
\noindent\hspace*{1.27cm}27. Yangi huquqiy mezonlar vakolatli hujjatda belgilangan sanadan qo‘llanadi. Ilgari qabul qilingan qarordagi xato o‘sha qarorga tatbiq etilgan huquqiy qoidalar bo‘yicha tekshiriladi; yangi mezon yoki keyingi davr parametri orqaga avtomatik tatbiq etilmaydi. Qayta hisoblashning o‘zi ilgari qonuniy tayinlangan to‘lovni qaytarib undirishga mustaqil asos bo‘lmaydi. To‘lovni o‘zgartirish, davom ettirish yoki undirish masalasi tegishli qonuniy asos va individual qaror bilan hal qilinadi. Qayta ko‘rib chiqish qaydi manba faktlari, avvalgi qoida ijrosi yoki huquqiy qoidaning o‘zi o‘zgarganini ajratadi; oldingi va tuzatilgan qismlar, tatbiq sanasi hamda individual qarorni o‘zgartirish asoslari saqlanadi. Ta’sirlangan ishlar ro‘yxati va vakolatli organga yuborish natijasi 17-banddagi nuqson qaydi bilan bog‘lanadi. To‘lovni o‘zgartirish uchun qonuniy asos bo‘lmasa, texnik tuzatish to‘lov topshirig‘iga aylanmaydi.\par
\noindent\hspace*{1.27cm}28. Agentlik saqlanadigan qaydlar toifalari, ularning qonuniy saqlash muddati va asosi, foydalanish vakolatlari, tuzatishlar tarixi, egasizlantirish yoki yo‘q qilish tartibini belgilaydi. Qaydlarni saqlash ushbu Tartib bo‘yicha cheksiz muddatli saqlashni anglatmaydi. Saqlash asosi yoki muddati tugagan ma’lumotlar qonunchilikdagi arxiv va shaxsga doir ma’lumotlar talablariga muvofiq yo‘q qilinadi yoxud qonuniy ruxsat etilgan holda egasizlantiriladi. Qonunchilik ma’lumotni yo‘q qilishni majburiy belgilagan hollarda egasizlantirish ushbu majburiyatning o‘rnini bosmaydi. Shikoyat yoki qonuniy nazorat uchun zarur qaydlar amaldagi saqlash talablariga muvofiq himoyalanadi. Saqlash toifalari kirish qiymati, hisoblash izi, nazorat bayonnomasi, tuzatish va murojaatni alohida qamraydi. Har biri maqsad, muddatning huquqiy asosi, mas’ul, kirish rollari va saqlashni tugatish usulini ko‘rsatadi. Himoyalangan yozuvga kirish va o‘zgartirish qayd etiladi; tekshiruvchi vakolati doirasida zarur hajmni oladi. Nazorat nusxasi qonuniy asos tugagach shaxsiy ma’lumot saqlash huquqini yaratmaydi.\par
\clearpage
\section*{Asosiy loyihaga kiritish uchun qisqa muqobil tahrir}
PF-193ning 14-bandi bo‘yicha hujjatga normativ kiritish uchun tavsiya etiladigan minimum. Quyidagi besh band yetti asosiy band va 28 bandli ilovaning to‘liq tahririni almashtiradi, unga qo‘shimcha majburiyat sifatida qo‘shilmaydi. Batafsil pasportlar, misollar va bayonnomalar ichki texnik hujjatlarga o‘tkaziladi. Muayyan talab yuqori hujjatlarga mos tegishli yuridik kuchdagi tatbiq norma bilan ta’minlangan bo‘lsa, ishlab chiquvchi shu normani saqlaydi va asosiy loyiha yoki uning asoslantirishida ko‘rsatadi; matnni takrorlash talab etilmaydi. Yakuniy raqamlash kiritiladigan joyga muvofiq belgilanadi.

1. O‘zbekiston Respublikasi Prezidenti huzuridagi Ijtimoiy himoya milliy agentligi (keyingi o‘rinlarda — Agentlik) O‘zbekiston Respublikasi Iqtisodiyot va moliya vazirligi (keyingi o‘rinlarda — Vazirlik) bilan birgalikda ko‘p o‘lchovli baholash uslubiyotini belgilangan tartibda tasdiqlashga tayyorlashda oilalarni toifalarga ajratish mezonlari va shartlarining huquqiy belgilanishini ta’minlasin. Bunda kiritish yoki rad etish shartini belgilaydigan qismda toifaga ta’sir qiluvchi chegaralar, tenglik, istisnolar ustuvorligi, birlik va davrlar, normativ daromad hisoblash, ruxsat etilgan kirish qiymatlari va yaxlitlash qamrab olinadi. Vakolatli hujjat, qo‘llash sanasi va saqlanadigan amaldagi qoidalarga havolalar ko‘rsatiladi. Dastur tavsifi va texnik qabul yetishmagan huquqiy shartni belgilamaydi yoki yuqori yuridik kuchdagi normani o‘zgartirmaydi. Normativ qoidalar belgilangan tartibda tasdiqlanadi va rasman e’lon qilinadi; idoraviy normativ hujjatga tatbiq etiladigan davlat ro‘yxati talabi bajariladi. Algoritmni vaqtincha o‘zgartirishga qonunchilikda bevosita berilgan vakolat o‘z predmeti doirasida saqlanadi.

2. Agentlik hisoblashning qonuniy qo‘llanadigan uslubiyotga muvofiqligini va ta’sirlangan qoidalarning birinchi qo‘llanishi yoki tegishli dasturiy o‘zgarishidan oldin mavjud chiqarilish jarayonida tekshirilishini ta’minlasin. Teng amaldagi tekshiruv takrorlanmasdan ishlatilsin. Natija huquqiy qoida, sana, sinalgan chiqarilish va mustaqil asoslangan kutilgan javobga bog‘lansin. Qarorga ta’sir qiladigan isbotlangan tafovut tuzatilib qayta tekshirilsin; boshqa dasturiy yig‘ilma faqat qaror sanasida qonuniy amal qiladigan qoidaga mosligi va ishlashga yaroqliligi tasdiqlansa qo‘llansin. Ta’sirlangan individual qarorlar vakolatli organga qonuniy qayta ko‘rish uchun yuborilsin. Majburiy texnik kelishuvlar, qonuniy muddatlar, boshlanish sanasi va yordam shartlari saqlansin.

3. Agentlik huquqiy ahamiyatli natija bilan birga hisoblashni qayta tiklashga yetarli bo‘lgan qo‘llangan huquqiy qoida, hisoblash talqini, kirish qiymatlari va davrlari, asosiy aylantirishlar, normativ hisoblash, istisno hamda toifa asosining qonuniy zarur qaydini saqlasin. Mavjud shaxsiy kabinet yoki Agentlikning “Inson” ijtimoiy xizmatlar markazi orqali arizachiga ochilishi mumkin bo‘lgan huquqiy va faktik asoslar, unga tegishli qiymat va davrlar hamda e’tiroz tartibi berilsin; natija haqidagi SMS izohni almashtirmaydi. Ma’lumot tuzatish va e’tirozni ko‘rishning qonuniy muddatlari saqlansin. Uchinchi shaxsning himoyalangan ma’lumoti ochilmasin; saqlash, kirish va saqlashni tugatish qonunchilikka muvofiq bajarilsin.

4. Baholashning huquqiy sharti o‘zgarishi vakolatli organga kiritilib, qonuniy sanadan qo‘llansin; dasturiy tuzatish o‘zgarmagan qonuniy qoidani bajarsin. Qonunchilik texnik nosozlik, almashish uzilishi yoki elektron ma’lumot to‘liq olinmaganda algoritmni vaqtincha o‘zgartirishni bevosita ruxsat etsa, vakolatli shaxs asos, ta’sirlangan manba va holatlar, vaqtinchalik qoida, boshlanish va tugatish shartlarini, tiklanishdan keyin zarur tekshiruv va qonuniy qayta ko‘rishni qayd etsin. Ushbu taklif mazkur vakolatni kengaytirmaydi. Yo‘q yoki zid huquqiy norma ta’sirlangan shart birinchi qo‘llanishidan oldin vakolatli qabul qiluvchi yoki sharhlovchi organ tomonidan hal qilinsin; sharh orqali o‘zgartirish kiritilmaydi. Texnik kutish rad etish, muddat uzaytirish, bekor qilingan mezonni qo‘llash yoki qarzdan ozod qilish asosini yaratmaydi.

5. Agentlik tekshiruv, shakl va qaydlarni mavjud vakolat hamda ishlab chiqish, kelishish va nazorat jarayonlarida tashkil etsin. Ichki texnik hujjat yangi arizachi majburiyati yoki huquqiy mezonni belgilamasin; teng spetsifikatsiya, sinov, qayd va kanallar qayta ishlatilsin. Qo‘shimcha xarajat majburiyatidan oldin Agentlik Vazirlik bilan birgalikda faqat qoplanmagan ish, ajratilgan xodim vaqti, shartnoma qamrovi va qonuniy moliyalashtirishni aniqlasin; arizani ishlash o‘zgarishi tegishli haqiqiy yuklama va vaqt daliliga asoslansin. Alohida qo‘lda xizmat, qo‘shimcha ariza bosqichi yoki yangi majburiy ochiq statistik hisobot joriy etilmasin.

\end{document}
